\begin{textsample}{POS Dim 1 – ba – Score 76.00 – ba/ba\_ef\_24.txt}  \label{ex:f1_pos_101}
8.5.1.
Geografia A intensificação dos processos globalizantes e a expansão do meio técnico científico informacional vêm provocando mudanças na sociedade e trazendo repercussões significativas para a educação no que \textbf{diz} respeito às suas práticas formativas.
Essa tendência vem estimulando a revisão de concepções acerca do currículo, \textbf{abordagens} epistemológicas, didáticas, \textbf{metodológicas} e políticas.
Por sua vez, o ensino da Geografia em sala de \textbf{aula} passou a ganhar novos desafios e \textbf{tantas} outras possibilidades de rearranjos pedagógicos nunca vistos antes em nossa história.
Todo o dinamismo mediado pelas \textbf{tecnologias} da informação e comunicação ( TICs ) vem promovendo novas formas de investigar, aprender, pensar e produzir o espaço de vivência cotidiana, \textbf{demandando} de professores e estudantes uma revisão dos \textbf{métodos} de produção, articulação e aplicação do conhecimento.
Com a expansão da circulação de pessoas, produtos, mercadorias e capital, a dinâmica social tem se tornado cada vez mais \textbf{complexa} e instável, reafirmando a posição de destaque que os saberes geográficos vêm assumindo nos processos formativos escolares.
Esse cenário promove novas oportunidades de interação entre as pessoas de diversas áreas do globo, proporcionando novas trocas diariamente.
Através dessas realidades, mediadas por redes \textbf{digitais}, os estudos sobre paisagem, região, território e lugar ganham novas proporções, em que conhecer lugares e pessoas se torna condição cada vez mais elementar para atender necessidades de sobrevivência e desenvolvimento da \textbf{humanidade}.
As noções de espaço e tempo, além das relações do ser humano com os meios sociais e naturais, vêm sendo alteradas intensamente.
A apreensão do conhecimento e a \textbf{compreensão} \textbf{crítica} do mundo, com o reconhecimento de suas semelhanças, diferenças, diversidades e desigualdades, tornam -se, cada dia mais, condições \textbf{indispensáveis} para a construção de uma sociedade mais justa e igualitária.
A necessidade de promover aprendizagens ativas na escola, que apontem para a formação de sujeitos \textbf{críticos}, autônomos, conscientes a respeito de si e do outro é cada vez maior.
Diante disso, ganham força as propostas que coadunam com ações integradoras por meio de práticas inter e transdisciplinares.
Novos olhares sobre a forma de aprender e ensinar comprometidos com a construção de \textbf{competências} vêm proporcionando cada vez maior desenvolvimento da \textbf{capacidade} humana de \textbf{mobilizar} habilidades, conteúdos \textbf{conceituais}, atitudinais e procedimentais com \textbf{vistas} a uma formação integral, fincadas em laços de solidariedade, cooperação, respeito à vida e ao meio ambiente.
Apesar disso, ainda persistem disparidades em relação ao acesso a bens básicos como saúde, segurança e educação.
De igual modo, intensifica -se a degradação da natureza nas micro e macroescalas, além da crescente violência e ainda permanência das desigualdades raciais e de gênero.
Essa realidade convida a escola a encarar essas metamorfoses impostas pela exclusão técnica e novas configurações das relações sociais.
Toda essa conjuntura vem sendo marcada por grandes instabilidades, flexibilidades e metamorfoses aceleradas no cenário político, \textbf{econômico} e social em diferentes escalas, com consideráveis repercussões \textbf{tanto} na Bahia como em todo o território brasileiro, \textbf{exigindo} o desenvolvimento de novos saberes específicos para lidar com esse espaço cada vez mais dinâmico, interativo, competitivo e excludente.
O respeito à dimensão humana dos indivíduos, considerando suas subjetividades, \textbf{capacidades} criativas, seu direito de ser, viver, sentir, expressar -se, respeitar a si e ao outro, passa pela apreensão cognitiva, \textbf{conceitual} e descritiva do espaço habitado, e por uma intervenção concreta que começa com práticas pedagógicas ativas e contextualizadas que partam de \textbf{demandas} locais.
O cotidiano dos estudantes \textbf{exige} novas \textbf{competências} e habilidades que extrapolem a \textbf{descrição} e \textbf{memorização} e \textbf{atinjam} as necessidades que se colocam na atualidade.
O desenvolvimento integral em suas dimensões física, cognitiva, afetiva, simbólica, ética, moral e social é fundamental para a promoção e \textbf{compreensão} do mundo no qual o \textbf{aluno} se insere.
A \textbf{capacidade} de abordar, discutir e intervir nas \textbf{demandas} da complexidade-mundo e as intervenções sociedade-natureza são \textbf{fundamentos} da Geografia.
A \textbf{Ciência} Geográfica contribui para formação cidadã, uma vez que reconhece e estimula o questionamento sobre a apropriação e transformação do arranjo sociedade-natureza, no âmbito local, regional, nacional ou global.
Dessa forma, contribui com uma aprendizagem continuada \textbf{baseada} em pensamento \textbf{crítico}, \textbf{reflexivo} e participativo.
O domínio do conhecimento geográfico em uma sociedade democrática é de fundamental importância para o exercício cidadão e \textbf{FIGURA} 9 As \textbf{Competências} Gerais, \textbf{Competências} da Área, Específicas da Geografia e as Habilidades Fonte : Elaboração dos redatores do Componente Curricular de Geografia ( 2018 ). 2.
Exercitar a curiosidade intelectual e recorrer à \textbf{abordagem} própria das \textbf{ciências}, incluindo a investigação, a reflexão, a \textbf{análise} \textbf{crítica}, a imaginação e a criatividade, para investigar causas, elaborar e testar hipóteses, formular e resolver \textbf{problemas} e criar soluções ( inclusive tecnológicas ) com base nos conhecimentos das diferentes áreas. formação das novas gerações.
Nas palavras de Cavalcanti³³ “ a Geografia consiste, portanto, num conjunto de conhecimentos constituídos da perspectiva da espacialidade.
Seu papel é \textbf{explicitar} a espacialidade das práticas sociais ”.
Para Filizola ( 2009, p. 24 ), o objetivo da Geografia Escolar é preparar um agir cotidiano ao circular com segurança no espaço, seja o de sua comunidade, seja o de outrem, demonstrando interesse pelo meio ambiente, tomando decisões, avaliando ações, assumindo posturas e atitudes \textbf{críticas} diante das mídias.³² As \textbf{competências} de Geografia a serem desenvolvidas no Ensino Fundamental estão voltadas para produção de um sujeito \textbf{reflexivo} e comprometido com a intervenção social através da tomada de consciência de si, do outro, de sua \textbf{localidade} e do mundo.
A partir da BNCC, são propostas sete \textbf{competências} básicas que transitam pelo direito de aprendizagem construído a partir da prática \textbf{reflexiva}, argumentação, aplicação e produção de conhecimentos sobre a vida coletiva, interação entre a sociedade e natureza, com uso dos conhecimentos \textbf{cartográficos} e técnicas de investigação geográfica. ³³ CAVALCANTI, Lana de Souza.
O ensino da Geografia na escola.
Campinas, SP : Papirus, 2012. p. 136.
Por meio dessa proposta, torna -se fundamental o \textbf{aprofundamento} do estudo do espaço capaz de promover a conexão entre diferentes temas em variadas escalas, envolvendo os princípios geográficos de analogia, conexão, diferenciação, distribuição, extensão, localização e ordem.
A formação do \textbf{raciocínio} socioespacial dos estudantes perpassa pela \textbf{compreensão} dos conceitos/categorias estruturantes da \textbf{Ciência} Geográfica, como : paisagem, lugar, território e região.
Esses conceitos são estratégicos para promover o desenvolvimento do pensamento espacial através do confronto entre a \textbf{ciência} e as experiências concretas do cotidiano.
A \textbf{Ciência} Geográfica \textbf{demanda}, constantemente, observação e \textbf{análise} do mundo construído cotidianamente, considerando a relação entre a sociedade e a natureza.
O estudo da distribuição dos elementos naturais e humanos, fenômenos de caráter geográfico, a \textbf{mobilização} de conteúdos para superação de \textbf{problemas} a partir da curiosidade, uso de \textbf{métodos} de investigação e elaboração de propostas coletivas mostram -se como atributos fundantes de conscientização humana, pois dão suporte ao exercício de reflexão e mudança de postura com intervenções concretas em consonância com a vida prática e real.
A contribuição da Geografia para o fortalecimento do currículo da Educação Fundamental requer um resgate à trajetória do pensamento geográfico, já que o ensino deste componente curricular foi fortemente \textbf{influenciado} pelas transformações na própria \textbf{ciência} em diferentes contextos históricos.
As contribuições do positivismo francês até as \textbf{abordagens} pós-modernas são fundamentais para a efetivação das políticas formativas significativas e atualizadas.
Nas últimas décadas, a Geografia passa a desenvolver novos trabalhos, privilegiando as dimensões subjetivas da relação humana com a natureza, considerando a cultura e o modo como se diversificam as percepções do espaço geográfico e as formas de sua configuração.
O espaço passa a ser visto por meio de sua singularidade, envolvendo outros saberes, principalmente a Sociologia, a Antropologia, as Ciências Política e Biológica, aprofundando sua identidade interdisciplinar.
O espaço então passa a ser, também, compreendido a partir das vivências dos grupos humanos e sua \textbf{correlação} entre valores, símbolos e comportamentos, como discutem Yi-Fu Tuan e Armand Frémont.
Essa tendência vem resgatar um conjunto de ideias, sentimentos e percepções que as pessoas têm do seu lugar de experiências, que tem o \textbf{potencial} de reforçar o compromisso cidadão das pessoas com as futuras gerações, como é expresso nesta proposta formativa em todo o Ensino Fundamental, desde o 1º até o 9º ano.
Essa aproximação entre a visão \textbf{crítica} e a percepção humanista \textbf{agregando} as vivências e o afeto entre os grupos sociais com o “ espaço vivido ”, considerando as dimensões simbólicas e estéticas dos indivíduos em seu cotidiano, envolve práticas capazes de ser contextualizadas nos diferentes territórios de identidade do nosso estado.
Ao entender o sentimento em relação aos lugares, torna -se inevitável questionar a descaracterização dos lugares em decorrência do processo evolutivo da globalização.
Tal processo \textbf{implica} a uniformização dos modos de vida e, consequentemente, dos espaços, como, por exemplo, a deterioração do meio ambiente em função do processo produtivo \textbf{capitalista}.
Em \textbf{vista} dessas questões, esta proposta curricular reforça a perspectiva de que todo ato educativo é político, pelo fato de ser manifestação de poder, referente aos saberes que todos possuem, interação e mediação destes com todos os integrantes do ato educativo.
As \textbf{abordagens} \textbf{críticas} são reforçadas, ao ser incluído o desenvolvimento da autonomia como uma das \textbf{competências} \textbf{centrais} a serem desenvolvidas na proposta educativa da escola.
O fortalecimento da “ autonomia ” nas práticas educativas é contemplado progressivamente em diferentes habilidades dentro de cada eixo temático e vinculado à ideia de participação social e política.
Reafirma -se a necessidade de privilegiar \textbf{competências} comprometidas com a transformação social a partir da \textbf{ótica} dos sujeitos da escola.
Partindo disso, a leitura de mundo e a promoção da autonomia do estudante no processo de aprendizagem geográfica passam a ser traduzidas como uma condição essencial para a \textbf{contextualização} da aprendizagem e diversificação das temáticas, objetos de conhecimento e conteúdos trabalhados pela escola.
Essa discussão está diretamente relacionada à própria construção da democracia, como o princípio inspirador do pensamento político-pedagógico e da gestão democrática.
A promoção da autonomia entre os sujeitos da educação requer o desenvolvimento de metodologias ativas de emancipação, sobretudo num contexto de instabilidade \textbf{econômica} e social marcada pela diminuição do emprego e aumento da violência como se vive neste \textbf{século}.
Não apenas pensar o mundo criticamente, mas também o desenvolvimento de atitudes responsáveis e éticas diante da realidade concreta são princípios fundantes para o ensino da Geografia.
As práticas pedagógicas neste componente são, por essência, interativas, dialogadas, privilegiando o questionamento, investigação e intervenção, partindo do local tendo em \textbf{vista} o global.
O ensino que se pretenda \textbf{relevante} deve ser comprometido com a superação de \textbf{problemas} sociais, \textbf{econômicos}, políticos, culturais e ambientais.
Assim, tornam -se necessárias as práticas que contemplem \textbf{aulas} mais atrativas, dinâmicas significativas e includentes.
Como, por exemplo, o uso de mapas temáticos, geoprocessamento, trabalho de campo, uso do GPS, dramatização, \textbf{entrevista}, videoaula, produção de \textbf{vídeo}, jornal falado, sala ambiente, confecção de painéis, criação de blog, grupos específicos nas redes sociais, leitura da paisagem, música, filmes, confecção de maquetes, fórum simulado etc.
Deve -se atentar para que essas práticas \textbf{agreguem} as múltiplas inteligências, sobretudo as voltadas para a educação inclusiva.
Um dos grandes desafios que enfrentamos na escola brasileira é a admissão dos \textbf{alunos} especiais, público-alvo da educação inclusiva.
Vivemos em uma sociedade altamente competitiva, excludente e preconceituosa.
A \textbf{devida} inserção do \textbf{aluno} especial “ está colocada como compromisso ético-político, que \textbf{implica} garantir a educação como direito de todos ” ( Prieto, 2006 ).
É dever da escola o reconhecimento de que todos têm a \textbf{capacidade} de apreender e ser respeitado em suas diferenças de sexo, orientação sexual, idade, classe social, etnia, língua, estado de saúde e deficiência.
Entretanto, as leis que regulamentam e asseguram esse direito ainda têm pouco avançado na prática para a inclusão educacional desses indivíduos no contexto do ensino regular.
Não obstante, os programas educacionais da escola regular ainda não contemplaram aqueles que necessitam de ações educativas integradoras, ou seja, para \textbf{alunos} com necessidades educacionais especiais conjuntamente com o ensino regular.
Em verdade, a flexibilidade dos conteúdos escolares, para a realidade do \textbf{aluno} especial no contexto escolar, ainda não contemplou a realidade de suas necessidades individuais ou coletivas para a real inserção na conjuntura escolar.
Desde séries iniciais, os \textbf{alunos} são \textbf{instigados} a desenvolver habilidades que venham a possibilitar uma \textbf{compreensão} do mundo a partir da perspectiva geográfica.
A possibilidade de leitura de mundo pode ser facilitada com a \textbf{introdução} de conhecimento específicos da \textbf{Ciência} Geográfica, sobretudo na \textbf{abordagem} da cartografia escolar.
Entretanto, esse entendimento dar -se-á mediante a metodologia concisa e voltada principalmente à vida cotidiana do \textbf{aluno}, a partir do conhecimento do próprio espaço corporal e da sua relação com o espaço vivido.
Conhecer o processo de mapeamento do espaço requer a criação de meios de representação mediante simbologia, aprender a ler informações geográficas contidas nas diversas representações \textbf{cartográficas}.
A leitura de um mapa \textbf{exige} questionamentos como : O que há em tal lugar?
Onde estão essas características?
Em que ordem?
Quanto?
É \textbf{imprescindível} o uso constante dessa representação gráfica em sala de \textbf{aula}.
Acrescente -se ainda que o ensino e a aprendizagem fazem parte de um mesmo processo, pois se \textbf{segue} um ao outro, e um sempre procede a outro.
Assim, não haverá ensino sem aprendizagem para o pleno desenvolvimento dos \textbf{alunos} com ou sem necessidades especiais.
É necessário o emprego de \textbf{métodos} de ensino que privilegiem suas potencialidades cognitivas, afetivas e motoras.
Isto só é possível, a partir de \textbf{métodos} educativos capazes de inseri -los no contexto do ensino regular.
A \textbf{simples} matrícula do \textbf{aluno} especial em séries regulares, sem o \textbf{devido} preparo didático-pedagógico, pode aumentar a resistência de alguns profissionais, como apontam os estudos de vários \textbf{autores} ( Bueno, 1998 ; Januzzi, 1992, 2004 ; Mazzotta, 1996 ).
Torna -se fundamental para a escola viabilizar ações que permitam cada vez mais a socialização desses indivíduos, pautadas no rigor científico.
No Brasil, ainda há poucos trabalhos em Geografia que abordem uma proposta \textbf{metodológica} para indivíduos especiais, entretanto temos renomados \textbf{pesquisadores} ( Almeida, 2010 ; Castrogiovanni, 2000 ; Callai, 2000 ; Martinelli, 2010 ; Siminelli, 2010 ; Paganelli, 2010 ) na linha de pesquisa sobre cartografia escolar que subsidiará o \textbf{aprofundamento} desta temática, pois serve de referência significativa no âmbito desta \textbf{abordagem}.
É importante ressaltar que, nesse contexto, a cartografia escolar pode \textbf{subsidiar} a \textbf{compreensão} do espaço em suas múltiplas \textbf{abordagens}, visto que essa área do conhecimento desvela a relação entre o espaço vivido e o espaço percebido.
Dessa forma, a cartografia escolar tem um lugar privilegiado na \textbf{compreensão}, pelo \textbf{aluno}, do espaço historicamente produzido.
A \textbf{interdisciplinaridade} é sem dúvida uma das bases epistemológicas da \textbf{Ciência} Geográfica, visto que dialoga e articula com outros saberes.
O conhecimento geográfico requer a construção de práticas significativas a partir de saberes que não podem ser fragmentados e descontextualizados.
História, Artes, Ciências da Natureza, Matemática e Linguagem apresentam textos, imagens, mapas, gráficos, ilustrações que exploram o conteúdo de forma interdisciplinar.
Assim, precisa -se assegurar o fortalecimento das relações entre as diferentes áreas do conhecimento para a real \textbf{compreensão} do todo.
Os múltiplos conhecimentos devem estar articulados com outros de caráter popular, filosófico e religioso a partir do contexto de vivência dos estudantes.
Devem -se considerar, nesse sentido, as especificidades, dialogando com o cotidiano dos \textbf{alunos}, das demais esferas da sua vida, como o lazer, as manifestações culturais, inclusive do trabalho.
É necessário esforço pedagógico e sistemático para considerar o mundo, a história, a cultura das populações quilombolas, indígenas, ribeirinhas, rurais, e as múltiplas modalidades de educação, como a Educação de Jovens e Adultos, a Educação Especial, Educação Prisional, dentre \textbf{tantos} outros presentes no Estado da Bahia.
Essa intenção formativa requer revisão da formação dos educadores, no seu perfil pedagógico e seu \textbf{posicionamento} ideológico além da visão do seu papel como profissional. \textbf{Somado} a isso, reafirma -se o compromisso com a qualidade no que \textbf{diz} respeito ao dever do Estado de garantir a modernização da estrutura física das escolas, adequações das instalações e atualização dos materiais que dão suporte às \textbf{aulas} práticas e lúdicas, como mapas, globos e bússolas, orçamento para \textbf{subsidiar} trabalho de campo e aquisição de equipamentos tecnológicos que deem suporte aos estudos e produções \textbf{cartográficas} por meio de mídias \textbf{digitais} e internet.
O ensino de Geografia, pautado no estudo da interação entre sociedade e natureza nas diferentes escalas espaciais, proporciona um campo \textbf{indispensável} para pensar nos \textbf{caminhos} para a organização, incorporação e \textbf{sistematização} do saber que os \textbf{alunos} constroem nas diversas esferas de suas vidas a partir do local onde vivem.
O estudo do meio favorece, de \textbf{maneira} especial, considerações sobre o que eles trazem das experiências do espaço e do tempo.
Para a garantia dos direitos de aprendizagens geográficas, \textbf{exige} -se esforço no sentido de promover atividades extraclasses, com exploração de múltiplos espaços e tempos dentro e fora da escola, com práticas inter e transdisciplinares, que visem sempre à \textbf{progressão} continuada dos estudos.
Considerar essas múltiplas realidades é um \textbf{caminho} importante para pensar em atividades educativas que respeitem o direito ao lazer e à diversão, muitas vezes reduzido a níveis muito baixos nas práticas cotidianas no que \textbf{diz} respeito ao ensino da Geografia.
Há de se considerar, também, a valorização da experiência.
Ela se \textbf{configura} através de tudo que passa entre os sujeitos, de tudo que acontece e que produz sentido para eles, inclusive o que os fazem viver ; ela é o que os \textbf{implica}, portanto os afeta, toca, \textbf{mobiliza} e também impõe e nos compromete.
Assim, a experiência pedagógica nunca os deixa indiferentes.
Adotar a noção de experiência para se promoverem práticas educativas \textbf{transformadoras} reside no interesse em valorizar os saberes multirreferenciados para além do que as tradições científicas e acadêmicas hierarquizadas instituem como válido.
Essa noção permite entender com maior \textbf{profundidade} como os indivíduos interpretam e organizam suas realidades e acabam construindo seus ordenamentos, ou seja, “ propõem e constroem investigações implicadas, engajadas ” ( MACEDO, 2015, p. 20 ), como se pensa numa aprendizagem espacial \textbf{contemporânea}.
Nessa lógica, a avaliação do avanço e desenvolvimento da aprendizagem e ensino deve ultrapassar as \textbf{provas} e testes escritos, historicamente privilegiados quando se valorizavam os conhecimentos \textbf{teóricos} e de forma pontual.
Já vem sendo amplamente discutida a importância de se diversificarem os instrumentos, os tempos e espaços de aferição, servindo de referência não apenas ao estudante, mas também ao professor.
Ela faz parte do processo formativo, daí a importância em ser desenvolvida processualmente, de forma dialogada, traçando estratégias para superação de dificuldades.
As estratégias \textbf{avaliativas} devem estar em sintonia com a prática desenvolvida e com os objetivos selecionados para cada ação pedagógica, a fim de prezar pela coerência e servir como estímulo ao avanço da investigação, dialogar no processo formativo e evitar exclusões e classificações desnecessárias.
É importante não \textbf{perder} de \textbf{vista} as finalidades da educação em geral e as especificidades da Geografia como uma \textbf{ciência} humana nem o objetivo do Ensino Fundamental.

% matched lemmas: abordagem, agregar, aluno, análise, aprofundamento, atingir, aula, autor, avaliativo, basear, caminho, capacidade, capitalista, cartográfico, central, ciência, competência, complexo, compreensão, conceitual, configurar, contemporâneo, contextualização, correlação, crítico, demanda, demandar, descrição, devido, digital, dizer, econômico, entrevista, exigir, explicitar, figurar, fundamento, humanidade, implicar, imprescindível, indispensável, influenciar, instigar, interdisciplinaridade, introdução, localidade, maneira, memorização, metodológico, mobilizar, mobilização, método, perder, pesquisador, posicionamento, potencial, problema, profundidade, progressão, prova, raciocínio, reflexivo, relevante, segar, simples, sistematização, somar, subsidiar, século, tanto, tecnologia, teórico, transformador, vista, vídeo, ótica
\end{textsample}
